%%=============================================================================
%% Methodologie
%%=============================================================================

\chapter{\IfLanguageName{dutch}{Methodologie}{Methodology}}%
\label{ch:methodologie}

%% TODO: In dit hoofstuk geef je een korte toelichting over hoe je te werk bent
%% gegaan. Verdeel je onderzoek in grote fasen, en licht in elke fase toe wat
%% de doelstelling was, welke deliverables daar uit gekomen zijn, en welke
%% onderzoeksmethoden je daarbij toegepast hebt. Verantwoord waarom je
%% op deze manier te werk gegaan bent.
%% 
%% Voorbeelden van zulke fasen zijn: literatuurstudie, opstellen van een
%% requirements-analyse, opstellen long-list (bij vergelijkende studie),
%% selectie van geschikte tools (bij vergelijkende studie, "short-list"),
%% opzetten testopstelling/PoC, uitvoeren testen en verzamelen
%% van resultaten, analyse van resultaten, ...
%%
%% !!!!! LET OP !!!!!
%%
%% Het is uitdrukkelijk NIET de bedoeling dat je het grootste deel van de corpus
%% van je bachelorproef in dit hoofstuk verwerkt! Dit hoofdstuk is eerder een
%% kort overzicht van je plan van aanpak.
%%
%% Maak voor elke fase (behalve het literatuuronderzoek) een NIEUW HOOFDSTUK aan
%% en geef het een gepaste titel.
Dit onderzoek volgt een ontwerpgericht en experimenteel onderzoeksdesign waarbij een proof-of-concept wordt ontwikkeld en geëvalueerd. Het doel is om na te gaan in welke mate AI-technieken ingezet kunnen worden voor automatische foutdetectie en feedbackgeneratie bij de kip aan barre.

Het onderzoeksproces wordt onderverdeeld in vijf grote fasen.

\section{Fase 1: Literatuurstudie}

In de eerste fase gebeurt een literatuurstudie uitgevoerd rond drie kerncomponenten: 

\begin{itemize}
    \item biomechanische kenmerken van de kip aan barre en veelvoorkomende fouten
    \item markerloze pose-estimatie in sportanalyse
    \item automatische foutdetectie en AI-gebaseerde feedbacksystemen
\end{itemize}

\textbf{Doelstelling:} theoretisch kader opbouwen en technische haalbaarheid inschatten.  

\textbf{Deliverable:} uitgewerkt literatuurhoofdstuk met onderbouwing van de gekozen aanpak.  

\textbf{Onderzoeksmethode:} systematische analyse van wetenschappelijke publicaties en bestaande toepassingen binnen sportanalyse.

\section{Fase 2: Specificatie van fouttypes en systeemvereisten}

\subsection{Specificatie van detecteerbare fouttypes}

Op basis van de literatuur en praktijkervaring binnen recreatieve gymnastiek worden concrete fouttypes gedefinieerd die automatisch detecteerbaar moeten zijn (zoals gebogen armen of onvoldoende heupsluiting).

Daarnaast worden systeemvereisten opgesteld die bepalen welke functionaliteiten het proof-of-concept minimaal moet bevatten. Hierbij wordt expliciet rekening gehouden met de praktische context van recreatieve gymnastiek, waarin gewerkt wordt met standaard videobeelden en beperkte technische infrastructuur.

Naast deze foutdefinities op basis van gewrichtshoeken wordt in de literatuur ook gewezen op het belang van kritieke tussenpunten (via-points) in de beweging \autocite{Yamasaki2010}.

Binnen dit onderzoek worden deze via-points initieel niet expliciet gemodelleerd in de eerste versie van het systeem. In plaats daarvan wordt gekozen voor een vereenvoudigde representatie op basis van fasen en gewrichtshoeken, om de technische haalbaarheid van automatische foutdetectie eerst te valideren.

De integratie van via-points als aanvullende analysemethode wordt voorzien als mogelijke uitbreiding in een iteratieve ontwikkelingsfase.

\subsection{Systeemvereisten}
\subsubsection{Functionele vereisten}

De functionele vereisten beschrijven welke functionaliteiten het systeem moet implementeren om de vooropgestelde doelstelling te bereiken. Voor dit proof-of-concept worden de volgende vereisten gedefinieerd:

\begin{enumerate}
    \item Het systeem moet een videofragment van een kip aan de barre kunnen inladen.
    \item Het systeem moet automatisch gewrichtsposities extraheren uit de video via markerloze pose-estimatie.
    \item Het systeem moet de beweging opdelen in relevante fasen (zwaai, trek, heupsluiting en steunfase).
    \item Het systeem moet vooraf gedefinieerde technische fouten detecteren op basis van pose-data:
    \begin{itemize}
        \item Gebogen armen tijdens de trekfase
        \item Onvoldoende heupsluiting
        \item Te vroege heupopening
        \item Onvoldoende schouderdruk in de steunfase
    \end{itemize}
    \item Het systeem moet de gedetecteerde fouten omzetten naar gestructureerde parameters.
    \item Het systeem moet op basis van deze parameters automatisch tekstuele feedback genereren via een Large Language Model.
    \item Het systeem moet de gegenereerde feedback tonen aan de gebruiker.
\end{enumerate}

\subsubsection{Niet-functionele vereisten}

Naast de functionele vereisten worden ook niet-functionele vereisten gedefinieerd. Deze beschrijven kwaliteitsaspecten waaraan het systeem moet voldoen en worden opgesteld volgens het SMART-principe.

\begin{table}[h]
    \centering
    \begin{tabular}{llll}
        \toprule
        NFR & Indicator & Meetvoorschrift & Norm \\
        \midrule
        Prestatie-efficiëntie & Snelheid & Tijd tussen input video en gegenereerde feedback & $\leq$ 10 s \\
        Functionele geschiktheid & Correctheid & Vergelijking met annotatie door expert & $\geq$ 70\% \\
        Bruikbaarheid & Begrijpelijkheid & Evaluatie door coach (Likert-schaal) & $\geq$ 3/5 \\
        Betrouwbaarheid & Foutbestendigheid & Testen met verschillende video’s & Geen crashes \\
        Overdraagbaarheid & Compatibiliteit & Ondersteunde videoformaten & Minstens MP4 \\
        \bottomrule
    \end{tabular}
    \caption{Niet-functionele vereisten}
\end{table}

\subsection{Prioritering van vereisten (MoSCoW)}

Om de scope van het proof-of-concept duidelijk af te bakenen, worden de vereisten geprioriteerd volgens de MoSCoW-methode.

\textbf{Must have}
\begin{itemize}
    \item Video-invoer
    \item Pose-estimatie
    \item Detectie van de vier gedefinieerde fouttypes
    \item Generatie van tekstuele feedback
\end{itemize}

\textbf{Should have}
\begin{itemize}
    \item Opsplitsing van de beweging in fasen
    \item Structurering van parameters
\end{itemize}

\textbf{Could have}
\begin{itemize}
    \item Visualisatie van het skelet op de video
    \item Vergelijking van meerdere pogingen
\end{itemize}

\textbf{Won't have}
\begin{itemize}
    \item Real-time analyse
    \item Mobiele applicatie
    \item Volledig gevalideerd coachesysteem
\end{itemize}

\textbf{Doelstelling:} afbakenen wat het systeem precies moet detecteren en genereren.  

\textbf{Deliverable:} lijst van fouttypes, systeemvereisten en bijhorende meetbare parameters.  

\textbf{Onderzoeksmethode:} conceptuele analyse en vertaling van biomechanische fouten naar kwantificeerbare bewegingskenmerken.

\section{Fase 3: Exploratief experiment met bestaande AI-systemen}

Voorafgaand aan de ontwikkeling van het proof-of-concept vindt een verkennend experiment plaats uitgevoerd met bestaande Large Language Models (LLMs), met als doel hun huidige capaciteiten voor gymnastiekanalyse te evalueren.

In dit experiment wordt eenzelfde kip-oefening aangeboden aan het model in drie verschillende representaties: (1) ruwe videobeelden, (2) pose-estimatie data in CSV-formaat, en (3) een grafische voorstelling van afgeleide gewrichtshoeken. Voor elk van deze inputvormen wordt een gelijkaardige prompt gebruikt waarbij het model wordt gevraagd om de beweging te analyseren, fasen te identificeren en technische feedback te formuleren.

\textbf{Doelstelling:} inzicht verkrijgen in de mogelijkheden en beperkingen van bestaande AI-systemen voor automatische bewegingsanalyse en feedbackgeneratie.

\textbf{Deliverable:} kwalitatieve analyse van de gegenereerde feedback per inputtype.

\textbf{Onderzoeksmethode:} exploratief experiment waarbij verschillende inputmodaliteiten systematisch worden vergeleken op basis van de gegenereerde output.

\section{Fase 4: Ontwikkeling van een proof-of-concept}

In deze fase wordt een proof-of-concept ontwikkeld waarin videobeelden worden verwerkt via markerloze pose-estimatie. Op basis van de geëxtraheerde keypoints worden relevante hoeken en bewegingsparameters berekend.

De videobeelden worden verwerkt met behulp van een Python-gebaseerde analysepipeline. Hiervoor wordt gebruikgemaakt van het MediaPipe Pose Landmarker model, dat markerloze pose-estimatie uitvoert op RGB-videoframes. Het gebruikte model detecteert per frame lichaamslandmarks die de positie van gewrichten beschrijven.

De gedetecteerde landmarks worden frame per frame geëxtraheerd en opgeslagen, waardoor een tijdreeks van lichaamsposities ontstaat. Op basis van deze gegevens worden bijkomende bewegingskenmerken berekend, zoals gewrichtshoeken. Deze kenmerken vormen de basis voor de verdere analyse van de beweging en de detectie van technische fouten in de kip-oefening.

De foutdetectie gebeurt via een rule-based benadering, waarbij vooraf gedefinieerde drempelwaarden bepalen of een fout aanwezig is. Gezien het exploratieve karakter van dit onderzoek en de beperkte beschikbaarheid van gelabelde trainingsdata wordt bewust gekozen voor een regelgebaseerde benadering in plaats van een volledig data-gedreven machine learning model. Deze aanpak verhoogt de interpreteerbaarheid van het systeem en laat toe om biomechanische criteria expliciet te definiëren.

Hoewel deze regelgebaseerde aanpak voornamelijk focust op drempelwaarden van gewrichtshoeken, sluit dit gedeeltelijk aan bij biomechanische inzichten uit de literatuur. 

Meer geavanceerde modellen, zoals beschreven door \textcite{Yamasaki2010}, suggereren dat bewegingen zoals de kip eerder bepaald worden door kritieke overgangsmomenten (via-points). Binnen dit proof-of-concept wordt echter gekozen voor een vereenvoudigde implementatie, waarbij deze dynamische aspecten impliciet benaderd worden via tijdsafhankelijke veranderingen in gewrichtshoeken.

Binnen de context van dit proof-of-concept vormt deze aanpak een geschikte eerste stap om de technische haalbaarheid van automatische foutdetectie te onderzoeken.

Gedetecteerde fouten worden vervolgens via een Large Language Model omgezet in gestructureerde tekstuele feedback.

Om de output van het systeem te structureren en consistent te maken, wordt een gestandaardiseerd dataformaat gedefinieerd in JSON-vorm. Dit formaat bevat informatie over de geïdentificeerde bewegingsfasen, berekende gewrichtshoeken en gedetecteerde technische fouten.

Deze gestructureerde representatie vormt de interface tussen de rule-based foutdetectie en de feedbackgeneratie via een Large Language Model. Om deze gegevens op een consistente manier door het taalmodel te laten interpreteren, wordt gebruikgemaakt van function calling en een vooraf gedefinieerd JSON-schema. Hierbij wordt het Large Language Model niet aangesproken met vrije tekstinput alleen, maar met een strikt gestructureerde datastructuur waarin per fouttype, bewegingsfase en relevante hoekwaarden expliciete velden zijn voorzien. Function calling laat toe het model te dwingen deze structuur correct te lezen en de output volgens een voorspelbaar formaat op te bouwen. Dit vermindert ambiguïteit in de interpretatie en verhoogt de reproduceerbaarheid van de gegenereerde feedback.

\textbf{Doelstelling:} aantonen dat automatische detectie en feedbackgeneratie technisch haalbaar zijn.  

\textbf{Deliverable:}  een proof-of-concept systeem dat videobeelden analyseert met behulp van pose-estimatie, bewegingskenmerken berekent en op basis daarvan automatische foutdetectie en feedbackgeneratie uitvoert.  

\textbf{Onderzoeksmethode:} experimentele implementatie en iteratieve ontwikkeling.

\section{Fase 5: Evaluatie}

Tot slot wordt het systeem geëvalueerd op twee niveaus: technische detectienauwkeurigheid en bruikbaarheid van de gegenereerde feedback.

De automatische foutdetectie wordt vergeleken met manuele beoordeling door een trainer. De feedbackoutput wordt kwalitatief beoordeeld op technische correctheid en toepasbaarheid binnen een trainingscontext.

Daarnaast wordt geëvalueerd in welke mate de huidige fase- en hoekgebaseerde benadering voldoende is om relevante technische fouten te detecteren, en of een uitbreiding naar een via-point gebaseerde analyse noodzakelijk is.

\textbf{Doelstelling:} inschatten in welke mate het systeem betrouwbaar en pedagogisch bruikbaar is.  

\textbf{Deliverable:} evaluatie-analyse en interpretatie van de resultaten.  

\textbf{Onderzoeksmethode:} vergelijkende analyse en kwalitatieve expertbeoordeling.

